%%%%%%%%%%%%%%%%%%%%%%%%%%%%%%%%%%%%%%%%%%%%%%%%%%%%%%%%%%%%%%%%%%%%%%%%%%%%%%%
%% TKS FORMAL MATHEMATICAL MANUAL v7.1 — ADVANCED FORMAL SEMANTICS
%% Continuity Bridge from v7.0
%% Extended Denotational Semantics with Domain Theory and Topology
%%%%%%%%%%%%%%%%%%%%%%%%%%%%%%%%%%%%%%%%%%%%%%%%%%%%%%%%%%%%%%%%%%%%%%%%%%%%%%%

\documentclass[11pt,twoside,openright]{book}

%% ============================================================================
%% PACKAGES (Preserved from v7.0 + v7.1 additions)
%% ============================================================================
\usepackage[utf8]{inputenc}
\usepackage[T1]{fontenc}
\usepackage{amsmath,amssymb,amsthm,amsfonts}
\usepackage{mathtools}
\usepackage{stmaryrd}           % For semantic brackets
\usepackage{bm}                 % Bold math
\usepackage{enumitem}
\usepackage{booktabs}
\usepackage{array}
\usepackage{longtable}
\usepackage{multirow}
\usepackage{graphicx}
\usepackage{tikz}
\usetikzlibrary{cd,arrows,positioning,calc,decorations.pathmorphing}
\usepackage{hyperref}
\usepackage{cleveref}
\usepackage{xcolor}
\usepackage{tcolorbox}
\tcbuselibrary{theorems,skins,breakable}
\usepackage{fancyhdr}
\usepackage{titlesec}
\usepackage{geometry}
\usepackage{listings}
\usepackage{multicol}

%% v7.0 PACKAGE ADDITIONS
\usepackage{braket}             % For quantum notation
\usepackage{tikz-cd}            % For commutative diagrams
\usepackage{bussproofs}         % For proof trees

%% v7.1 PACKAGE ADDITIONS
\usepackage{mathrsfs}           % For \mathscr (Scott topology)

\geometry{
  paper=letterpaper,
  inner=1.2in,
  outer=1in,
  top=1in,
  bottom=1in
}

%% ============================================================================
%% THEOREM ENVIRONMENTS (Preserved from v7.0)
%% ============================================================================
\theoremstyle{definition}
\newtheorem{definition}{Definition}[chapter]
\newtheorem{axiom}[definition]{Axiom}

\theoremstyle{plain}
\newtheorem{theorem}[definition]{Theorem}
\newtheorem{lemma}[definition]{Lemma}
\newtheorem{proposition}[definition]{Proposition}
\newtheorem{corollary}[definition]{Corollary}

\theoremstyle{remark}
\newtheorem{remark}[definition]{Remark}
\newtheorem{example}[definition]{Example}
\newtheorem{notation}[definition]{Notation}

%% ============================================================================
%% CUSTOM COMMANDS (Preserved from v7.0)
%% ============================================================================
% Semantic brackets
\newcommand{\sem}[1]{\llbracket #1 \rrbracket}

% TKS-specific symbols
\newcommand{\Worlds}{\mathcal{W}}
\newcommand{\Ideas}{\mathcal{I}}
\newcommand{\Elements}{\mathcal{E}}
\newcommand{\Noetics}{\mathcal{N}}
\newcommand{\Foundations}{\mathcal{F}}
\newcommand{\Acquisitions}{\mathcal{A}}
\newcommand{\Noetica}{\mathbf{Noetica}}
\newcommand{\Acq}{\mathbf{Acq}}

% Noetic operators
\newcommand{\noe}[1]{\nu_{#1}}

% Tootra operations
\newcommand{\Tadd}{\oplus_T}
\newcommand{\Tsub}{\ominus_T}
\newcommand{\Tmul}{\otimes_T}
\newcommand{\Tdiv}{\oslash_T}

% World association
\newcommand{\Wadd}{\oplus_W}
\newcommand{\Wsub}{\ominus_W}

% Type judgments
\newcommand{\types}{\vdash}
\newcommand{\typmark}{:}

% RPM Monad
\newcommand{\RPM}{\mathfrak{R}}
\newcommand{\rpmret}{\mathsf{return}}
\newcommand{\rpmbind}{\mathbin{\gg\!=}}
\newcommand{\Kleisli}{\mathcal{K}}

% Fractal notation
\newcommand{\Frac}{\Phi}
\newcommand{\FracSet}{\mathfrak{F}}
\newcommand{\bigstep}{\Downarrow}
\newcommand{\smallstep}{\rightarrow}
\newcommand{\Failed}{\mathsf{Failed}}

% Quantum notation (v7.0)
\newcommand{\Hspace}{\mathcal{H}}
\newcommand{\qket}[1]{|#1\rangle}
\newcommand{\qbra}[1]{\langle #1|}
\newcommand{\qbraket}[2]{\langle #1 | #2 \rangle}
\newcommand{\qop}[1]{\hat{#1}}

%% v7.1 NEW COMMANDS
% Domain theory
\newcommand{\Scott}{\mathscr{S}}           % Scott topology
\newcommand{\Dcpo}{\mathbf{Dcpo}}          % Category of dcpos
\newcommand{\Cont}{\mathbf{Cont}}          % Continuous functions
\newcommand{\sqleq}{\sqsubseteq}           % Information ordering
\newcommand{\bigsqcup}{\bigsqcup}          % Directed supremum
\newcommand{\fix}{\mathsf{fix}}            % Fixed point operator
\newcommand{\bot}{\perp}                   % Bottom element

% Natural transformations
\newcommand{\NatTrans}[2]{\eta_{#1 \to #2}}
\newcommand{\WorldFunctor}[1]{\mathfrak{W}_{#1}}

% Version markers
\newcommand{\vnew}{\textbf{v7.0 New}}
\newcommand{\vexp}{\textbf{v7.0 Expanded}}
\newcommand{\vnewone}{\textbf{v7.1 New}}
\newcommand{\vexpone}{\textbf{v7.1 Expanded}}

%% ============================================================================
%% DOCUMENT METADATA
%% ============================================================================
\title{\Huge\textbf{The TKS Formal Mathematical Manual}\\[0.5em]
       \Large Version 7.1 — Advanced Formal Semantics\\[0.3em]
       \large Continuity Bridge from v7.0}
\author{TKS Formalization Project}
\date{\today}

\begin{document}

\frontmatter
\maketitle

%% ============================================================================
%% PREFACE: v7.0 → v7.1 CONTINUITY BRIDGE
%% ============================================================================

\chapter*{v7.0 $\to$ v7.1 Continuity Bridge}
\addcontentsline{toc}{chapter}{v7.0 → v7.1 Continuity Bridge}

\begin{tcolorbox}[colback=yellow!10!white,colframe=orange!75!black,title=\textbf{Canonical Inheritance Declaration}]
This document, TKS v7.1, is a \textbf{strict superset} of TKS v7.0. All definitions, theorems, axioms, and proofs from v7.0 remain valid and unchanged. Version 7.1 introduces \textit{extensions only}---no modifications to existing canon.
\end{tcolorbox}

\section*{Inheritance from v7.0}

The following components from v7.0 are inherited verbatim:

\begin{enumerate}
  \item \textbf{Canonical Ontology}: 40 Elements, 4 Worlds, 10 Noetics, 7 Foundations, 28 Sub-Foundations, 22 Acquisitions
  \item \textbf{Noetic Operator Algebra}: Complete 10×10 composition table and idempotent structure
  \item \textbf{Tootra Arithmetic}: All four operations with multiset semantics
  \item \textbf{Category Noetica}: Objects $\{D_1, \ldots, D_7\}$, morphisms as noetic compositions
  \item \textbf{ACBE Functor}: $\mathrm{ACBE}: \Noetica_{|A} \to \Noetica_{|D}$
  \item \textbf{RPM Monad Theory}: Complete Kleisli category with effect system
  \item \textbf{Noetic Fractal Calculus}: Fractal sequences, composition, eigenfractals
  \item \textbf{Dynamic Semantics}: Evaluation environments, big-step and small-step rules
  \item \textbf{Type System}: Complete with subtyping and polymorphism
  \item \textbf{Concurrency Model}: Parallel composition and synchronization primitives
  \item \textbf{Quantum Noetic Algebra}: Hilbert space representation and measurement
\end{enumerate}

\section*{v7.1 Extensions}

Version 7.1 introduces the following new developments:

\begin{enumerate}
  \item \textbf{Natural Transformations Between Worlds} (Chapter~\ref{ch:nat-trans})
  \begin{itemize}
    \item Category-theoretic morphisms $\eta_{W \to W'} : \WorldFunctor{W} \Rightarrow \WorldFunctor{W'}$
    \item Coherence conditions for world transitions
    \item ACBE as a natural transformation
  \end{itemize}

  \item \textbf{Domain-Theoretic Semantics} (Chapter~\ref{ch:domain-theory})
  \begin{itemize}
    \item TKS domains as directed-complete partial orders (dcpos)
    \item Scott topology on semantic domains
    \item Continuous function semantics
    \item Fixed-point theorems for recursive definitions
  \end{itemize}

  \item \textbf{Extended Denotational Semantics} (Chapter~\ref{ch:extended-denotational})
  \begin{itemize}
    \item Full abstraction theorem
    \item Logical relations for semantic equivalence
    \item Metric semantics for convergence properties
  \end{itemize}

  \item \textbf{Topological Structure of Noetic Space} (Chapter~\ref{ch:noetic-topology})
  \begin{itemize}
    \item Noetic manifolds and continuous transformations
    \item Sheaf-theoretic interpretation of world localization
    \item Cohomological invariants of fractal structures
  \end{itemize}
\end{enumerate}

\section*{Notation Changes}

\begin{center}
\begin{tabular}{lll}
\toprule
\textbf{Symbol} & \textbf{Meaning} & \textbf{Introduced} \\
\midrule
$\Scott$ & Scott topology & v7.1 \\
$\Dcpo$ & Category of dcpos & v7.1 \\
$\sqleq$ & Information ordering & v7.1 \\
$\fix$ & Fixed-point operator & v7.1 \\
$\NatTrans{W}{W'}$ & Natural transformation & v7.1 \\
$\WorldFunctor{W}$ & World functor & v7.1 \\
\bottomrule
\end{tabular}
\end{center}

\section*{Backwards Compatibility}

All v7.0 programs and proofs are valid in v7.1. The new extensions provide:
\begin{itemize}
  \item Deeper semantic foundations (domain theory)
  \item Categorical understanding of world transitions
  \item Topological tools for analyzing noetic structures
  \item Stronger metatheoretic results (full abstraction)
\end{itemize}

No existing theorem is invalidated; new theorems extend the theory.

\tableofcontents

\mainmatter

%% ============================================================================
%% CANONICAL REFERENCE: 28 SUB-FOUNDATIONS (Inherited from v7.0)
%% ============================================================================
\part{Canonical Reference}

\chapter{The 28 Sub-Foundations}
\label{ch:28-subfoundations}

\begin{tcolorbox}[colback=gray!10!white,colframe=gray!75!black,title=\textbf{Canonical Reference (Inherited)}]
This chapter documents the 28 Sub-Foundations, inherited unchanged from v7.0/v6.1. The sub-foundations are the 7 Foundations projected across the 4 Worlds.
\end{tcolorbox}

\begin{definition}[Sub-Foundation Space]
\label{def:v71-subfoundation-space}
The \textbf{Sub-Foundation space} is the product:
\[
\Foundations^* = \Foundations \times \Worlds = \{F_1, \ldots, F_7\} \times \{A, B, C, D\}
\]
with $|\Foundations^*| = 7 \times 4 = 28$ sub-foundations.
\end{definition}

\begin{table}[ht]
\centering
\caption{The 28 Sub-Foundations}
\label{tab:v71-28-subfoundations}
\begin{tabular}{c|cccc}
\toprule
\textbf{Foundation} & \textbf{A (Spiritual)} & \textbf{B (Mental)} & \textbf{C (Emotional)} & \textbf{D (Physical)} \\
\midrule
$F_1$ (Unity) & $1_A$ & $1_B$ & $1_C$ & $1_D$ \\
$F_2$ (Wisdom) & $2_A$ & $2_B$ & $2_C$ & $2_D$ \\
$F_3$ (Life) & $3_A$ & $3_B$ & $3_C$ & $3_D$ \\
$F_4$ (Companionship) & $4_A$ & $4_B$ & $4_C$ & $4_D$ \\
$F_5$ (Power) & $5_A$ & $5_B$ & $5_C$ & $5_D$ \\
$F_6$ (Material) & $6_A$ & $6_B$ & $6_C$ & $6_D$ \\
$F_7$ (Lust) & $7_A$ & $7_B$ & $7_C$ & $7_D$ \\
\bottomrule
\end{tabular}
\end{table}

\begin{proposition}[Sub-Foundation Lattice Structure]
\label{prop:v71-subfoundation-lattice}
$(\Foundations^*, \sqcup, \sqcap)$ forms a bounded distributive lattice with:
\begin{enumerate}[label=(\roman*)]
  \item Product decomposition: $\Foundations^* \cong \prod_{w \in \Worlds} \Foundations_w$
  \item ACBE-compatible ordering: $i_A \leq i_B \leq i_C \leq i_D$ for each foundation $i$
  \item World projection: $\pi_w : \Foundations^* \to \Foundations_w$ extracts the $w$-level sub-foundations
\end{enumerate}
\end{proposition}

\begin{theorem}[Sub-Foundation Natural Transformation]
\label{thm:subfoundation-nat-trans}
The ACBE cascade induces a natural transformation on sub-foundations:
\[
\alpha : \Foundations^*_A \Rightarrow \Foundations^*_D
\]
with components $\alpha_i : i_A \to i_D$ for each foundation $F_i$, satisfying:
\[
\alpha_i = \mathsf{desc}_D \circ \mathsf{desc}_C \circ \mathsf{desc}_B
\]
where $\mathsf{desc}_w : \Foundations^*_{w'} \to \Foundations^*_w$ is the world descent morphism.
\end{theorem}

%% ============================================================================
%% PART II: NATURAL TRANSFORMATIONS BETWEEN WORLDS
%% ============================================================================
\part{Categorical Extensions}

\chapter{Natural Transformations Between Worlds}
\label{ch:nat-trans}

\begin{tcolorbox}[colback=blue!5!white,colframe=blue!75!black,title=\vnewone]
This chapter develops the category-theoretic structure of world transitions, formalizing the ACBE cascade and general inter-world morphisms as natural transformations.
\end{tcolorbox}

\section{World Functors}

\begin{definition}[World Category]
\label{def:world-cat}
Let $\mathbf{World}$ be the category where:
\begin{itemize}
  \item Objects: $\{A, B, C, D\}$ (the four worlds)
  \item Morphisms:
  \begin{align}
    \mathrm{Hom}(A, B) &= \{\iota_{AB}\} \quad \text{(descent from Spiritual to Mental)} \\
    \mathrm{Hom}(B, C) &= \{\iota_{BC}\} \quad \text{(descent from Mental to Astral)} \\
    \mathrm{Hom}(C, D) &= \{\iota_{CD}\} \quad \text{(descent from Astral to Physical)} \\
    \mathrm{Hom}(W, W) &= \{\mathrm{id}_W\} \quad \text{(identity)}
  \end{align}
  \item Composition: Sequential descent (non-reversible)
\end{itemize}
\end{definition}

\begin{remark}
The category $\mathbf{World}$ is a total order category (poset): $A > B > C > D$ with morphisms representing descent. There are no ascending morphisms—this reflects the hermetic principle that manifestation flows downward from Spirit to Matter.
\end{remark}

\begin{definition}[World Functor]
\label{def:world-functor}
For each world $W \in \{A, B, C, D\}$, define the \textbf{world functor}:
\[
\WorldFunctor{W} : \Noetica \to \mathbf{Set}
\]
by:
\begin{align}
\WorldFunctor{W}(D_j) &= \{Wn \mid n \in \{1, \ldots, 10\}, \mathrm{domain}(Wn) = D_j\} \\
\WorldFunctor{W}(f: D_i \to D_j) &= f \circ (-) : \WorldFunctor{W}(D_i) \to \WorldFunctor{W}(D_j)
\end{align}
\end{definition}

\begin{proposition}[Functor Laws]
\label{prop:world-functor-laws}
Each $\WorldFunctor{W}$ satisfies:
\begin{enumerate}
  \item $\WorldFunctor{W}(\mathrm{id}_{D_j}) = \mathrm{id}_{\WorldFunctor{W}(D_j)}$
  \item $\WorldFunctor{W}(g \circ f) = \WorldFunctor{W}(g) \circ \WorldFunctor{W}(f)$
\end{enumerate}
\end{proposition}

\begin{proof}
(1) For any $Wn \in \WorldFunctor{W}(D_j)$:
\[
\WorldFunctor{W}(\mathrm{id}_{D_j})(Wn) = \mathrm{id}_{D_j} \circ Wn = Wn = \mathrm{id}(Wn)
\]

(2) For $f: D_i \to D_j$ and $g: D_j \to D_k$:
\[
\WorldFunctor{W}(g \circ f)(Wn) = (g \circ f)(Wn) = g(f(Wn)) = \WorldFunctor{W}(g)(\WorldFunctor{W}(f)(Wn))
\]
\end{proof}

\section{Inter-World Natural Transformations}

\begin{definition}[World Descent Transformation]
\label{def:descent-trans}
For worlds $W > W'$ (with $W$ above $W'$ in the hierarchy), the \textbf{descent transformation}:
\[
\NatTrans{W}{W'} : \WorldFunctor{W} \Rightarrow \WorldFunctor{W'}
\]
is defined componentwise for each domain $D_j$:
\[
(\NatTrans{W}{W'})_{D_j} : \WorldFunctor{W}(D_j) \to \WorldFunctor{W'}(D_j)
\]
\[
(\NatTrans{W}{W'})_{D_j}(Wn) = W'n
\]
\end{definition}

\begin{theorem}[Naturality of Descent]
\label{thm:descent-naturality}
The descent transformation $\NatTrans{W}{W'}$ is natural: for any morphism $f: D_i \to D_j$ in $\Noetica$, the following diagram commutes:
\[
\begin{tikzcd}
\WorldFunctor{W}(D_i) \arrow[r, "\WorldFunctor{W}(f)"] \arrow[d, "(\NatTrans{W}{W'})_{D_i}"']
  & \WorldFunctor{W}(D_j) \arrow[d, "(\NatTrans{W}{W'})_{D_j}"] \\
\WorldFunctor{W'}(D_i) \arrow[r, "\WorldFunctor{W'}(f)"']
  & \WorldFunctor{W'}(D_j)
\end{tikzcd}
\]
\end{theorem}

\begin{proof}
For any $Wn \in \WorldFunctor{W}(D_i)$:

\textbf{Clockwise path:}
\begin{align}
(\NatTrans{W}{W'})_{D_j}(\WorldFunctor{W}(f)(Wn)) &= (\NatTrans{W}{W'})_{D_j}(f(Wn)) \\
&= f(W'n) \quad \text{(applying descent to result)}
\end{align}

\textbf{Counterclockwise path:}
\begin{align}
\WorldFunctor{W'}(f)((\NatTrans{W}{W'})_{D_i}(Wn)) &= \WorldFunctor{W'}(f)(W'n) \\
&= f(W'n)
\end{align}

Both paths yield $f(W'n)$, establishing naturality.
\end{proof}

\section{ACBE as Composite Natural Transformation}

\begin{theorem}[ACBE Naturality]
\label{thm:acbe-natural}
The ACBE functor is the composite natural transformation:
\[
\mathrm{ACBE} = \NatTrans{C}{D} \circ \NatTrans{B}{C} \circ \NatTrans{A}{B} : \WorldFunctor{A} \Rightarrow \WorldFunctor{D}
\]
\end{theorem}

\begin{proof}
For any $An \in \WorldFunctor{A}(D_j)$:
\begin{align}
\mathrm{ACBE}_{D_j}(An) &= ((\NatTrans{C}{D})_{D_j} \circ (\NatTrans{B}{C})_{D_j} \circ (\NatTrans{A}{B})_{D_j})(An) \\
&= (\NatTrans{C}{D})_{D_j}((\NatTrans{B}{C})_{D_j}(Bn)) \\
&= (\NatTrans{C}{D})_{D_j}(Cn) \\
&= Dn
\end{align}

This matches Definition 6.1 of ACBE from v7.0.
\end{proof}

\begin{corollary}[ACBE Coherence]
\label{cor:acbe-coherence}
For any noetic morphism $\noe{k} : D_i \to D_j$:
\[
\mathrm{ACBE}_{D_j} \circ \WorldFunctor{A}(\noe{k}) = \WorldFunctor{D}(\noe{k}) \circ \mathrm{ACBE}_{D_i}
\]
\end{corollary}

\section{Vertical Composition and 2-Category Structure}

\begin{definition}[World 2-Category]
\label{def:world-2-cat}
Define the 2-category $\mathbf{World}_2$ where:
\begin{itemize}
  \item 0-cells: Worlds $\{A, B, C, D\}$
  \item 1-cells: World functors $\WorldFunctor{W}$
  \item 2-cells: Natural transformations $\NatTrans{W}{W'}$
\end{itemize}
\end{definition}

\begin{proposition}[Vertical Composition]
\label{prop:vertical-comp}
Natural transformations compose vertically:
\[
\NatTrans{W}{W''} = \NatTrans{W'}{W''} \bullet \NatTrans{W}{W'} \quad \text{for } W > W' > W''
\]
\end{proposition}

\begin{theorem}[Strict 2-Functor]
\label{thm:strict-2-functor}
The assignment $W \mapsto \WorldFunctor{W}$ extends to a strict 2-functor:
\[
\mathfrak{W} : \mathbf{World}_2 \to \mathbf{Cat}
\]
preserving vertical and horizontal composition strictly.
\end{theorem}

%% ============================================================================
%% PART II: DOMAIN-THEORETIC SEMANTICS
%% ============================================================================
\part{Extended Formal Semantics}

\chapter{Domain-Theoretic Foundations}
\label{ch:domain-theory}

\begin{tcolorbox}[colback=blue!5!white,colframe=blue!75!black,title=\vnewone]
This chapter develops the domain-theoretic foundations for TKS semantics, providing a rigorous basis for recursive definitions, fixed points, and semantic continuity.
\end{tcolorbox}

\section{Directed-Complete Partial Orders}

\begin{definition}[Partial Order on TKS Values]
\label{def:tks-po}
Define the \textbf{information ordering} $\sqleq$ on TKS values:
\begin{enumerate}
  \item $\bot \sqleq v$ for all values $v$ (bottom is least)
  \item For elements: $Wn \sqleq Wn$ (reflexive, no proper refinement)
  \item For RPM values: $\Failed(p) \sqleq \Failed(p)$ and $\rpmret\; v \sqleq \rpmret\; v'$ iff $v \sqleq v'$
  \item For fractals: $\Frac \sqleq \Frac'$ iff $\Frac$ is a prefix of $\Frac'$
  \item For functions: $f \sqleq g$ iff $\forall x.\, f(x) \sqleq g(x)$
\end{enumerate}
\end{definition}

\begin{definition}[Directed Set]
\label{def:directed}
A subset $S \subseteq D$ of a poset $(D, \sqleq)$ is \textbf{directed} if:
\begin{enumerate}
  \item $S \neq \varnothing$
  \item For all $x, y \in S$, there exists $z \in S$ with $x \sqleq z$ and $y \sqleq z$
\end{enumerate}
\end{definition}

\begin{definition}[Dcpo]
\label{def:dcpo}
A poset $(D, \sqleq)$ is a \textbf{directed-complete partial order (dcpo)} if every directed subset $S \subseteq D$ has a supremum $\bigsqcup S \in D$.
\end{definition}

\begin{theorem}[TKS Semantic Domains are Dcpos]
\label{thm:tks-dcpo}
The following TKS semantic domains are dcpos:
\begin{enumerate}
  \item $\sem{\mathsf{Element}}_\bot = \{Wn \mid W \in \Worlds, n \in \{1,\ldots,10\}\} \cup \{\bot\}$
  \item $\sem{\mathsf{RPM}[\tau]}_\bot$ for any type $\tau$
  \item $\sem{\tau \to \sigma}$ for dcpo types $\tau, \sigma$
  \item $\sem{\mathsf{Frac}}_\bot$ (fractals with prefix ordering)
\end{enumerate}
\end{theorem}

\begin{proof}
(1) $\sem{\mathsf{Element}}_\bot$ is a flat domain: all non-$\bot$ elements are incomparable, so every directed set is either $\{\bot\}$ or contains a unique non-$\bot$ element.

(2) $\sem{\mathsf{RPM}[\tau]}_\bot$ inherits dcpo structure from $\tau$: directed sets of successful computations converge to successful computations (or remain at $\bot$ if all fail to terminate).

(3) Function spaces: For directed $\{f_i\}_{i \in I}$, define $(\bigsqcup_i f_i)(x) = \bigsqcup_i f_i(x)$. This is well-defined since each $\{f_i(x)\}$ is directed.

(4) Fractals with prefix ordering form an $\omega$-algebraic domain with compact elements being finite fractals.
\end{proof}

\section{Scott Topology}

\begin{definition}[Scott Open Sets]
\label{def:scott-open}
For a dcpo $(D, \sqleq)$, a subset $U \subseteq D$ is \textbf{Scott open} if:
\begin{enumerate}
  \item $U$ is upward closed: $x \in U$ and $x \sqleq y$ implies $y \in U$
  \item $U$ is inaccessible by directed suprema: for any directed $S$ with $\bigsqcup S \in U$, there exists $s \in S$ with $s \in U$
\end{enumerate}
The collection of Scott open sets forms the \textbf{Scott topology} $\Scott(D)$.
\end{definition}

\begin{proposition}[Scott Topology on TKS Domains]
\label{prop:scott-tks}
For the flat domain $\sem{\mathsf{Element}}_\bot$:
\[
\Scott(\sem{\mathsf{Element}}_\bot) = \{\varnothing\} \cup \{U \subseteq \sem{\mathsf{Element}}_\bot \mid \bot \notin U\} \cup \{\sem{\mathsf{Element}}_\bot\}
\]
The non-trivial open sets are exactly the non-empty upward-closed sets not containing $\bot$.
\end{proposition}

\begin{definition}[Scott Continuous Function]
\label{def:scott-continuous}
A function $f : D \to E$ between dcpos is \textbf{Scott continuous} if:
\begin{enumerate}
  \item $f$ is monotone: $x \sqleq y \implies f(x) \sqleq f(y)$
  \item $f$ preserves directed suprema: $f(\bigsqcup S) = \bigsqcup f(S)$ for directed $S$
\end{enumerate}
\end{definition}

\begin{theorem}[Noetic Operators are Scott Continuous]
\label{thm:noetic-scott}
Each noetic operator $\mathfrak{n}_k : \sem{\mathsf{Domain}}_\bot \to \sem{\mathsf{Domain}}_\bot$ is Scott continuous.
\end{theorem}

\begin{proof}
Monotonicity: Since $\sem{\mathsf{Domain}}_\bot$ is flat, monotonicity reduces to:
\begin{itemize}
  \item $\mathfrak{n}_k(\bot) = \bot$ (strictness)
  \item $\mathfrak{n}_k(D_i) \sqleq \mathfrak{n}_k(D_i)$ (reflexivity)
\end{itemize}

Directed suprema: For directed $S \subseteq \sem{\mathsf{Domain}}_\bot$:
\begin{itemize}
  \item If $S = \{\bot\}$, then $\bigsqcup S = \bot$ and $\mathfrak{n}_k(\bot) = \bot = \bigsqcup \mathfrak{n}_k(S)$
  \item If $D_i \in S$ for some $i$, then $\bigsqcup S = D_i$ (flatness) and preservation follows
\end{itemize}
\end{proof}

\section{Fixed-Point Theorems}

\begin{theorem}[Kleene Fixed-Point Theorem for TKS]
\label{thm:kleene-tks}
For any Scott continuous $f : D \to D$ on a pointed dcpo $(D, \bot, \sqleq)$:
\[
\fix(f) = \bigsqcup_{n \geq 0} f^n(\bot)
\]
exists and is the least fixed point of $f$.
\end{theorem}

\begin{proof}
The sequence $\bot \sqleq f(\bot) \sqleq f^2(\bot) \sqleq \cdots$ is an $\omega$-chain (hence directed). By dcpo property, $\bigsqcup_n f^n(\bot)$ exists.

Let $x^* = \bigsqcup_n f^n(\bot)$. Then:
\begin{align}
f(x^*) &= f(\bigsqcup_n f^n(\bot)) \\
&= \bigsqcup_n f^{n+1}(\bot) \quad \text{(continuity)} \\
&= \bigsqcup_{m \geq 1} f^m(\bot) \\
&= x^* \quad \text{(since } f^0(\bot) = \bot \text{ is least)}
\end{align}

For leastness: if $f(y) = y$, then by induction $f^n(\bot) \sqleq y$ for all $n$, so $x^* = \bigsqcup_n f^n(\bot) \sqleq y$.
\end{proof}

\begin{corollary}[Recursive Fractal Semantics]
\label{cor:recursive-fractal}
For a recursive fractal definition $\Frac = \langle k_1 : k_2 : \cdots \rangle(\Frac)$, the semantics is:
\[
\sem{\Frac} = \fix(\lambda x. \mathfrak{n}_{k_1}(\mathfrak{n}_{k_2}(\cdots(x))))
\]
\end{corollary}

\section{Continuous Lattice Structure}

\begin{definition}[Way-Below Relation]
\label{def:way-below}
For elements $x, y$ in a dcpo $D$, we say $x$ is \textbf{way below} $y$, written $x \ll y$, if:
\[
\text{For every directed } S \text{ with } y \sqleq \bigsqcup S, \text{ there exists } s \in S \text{ with } x \sqleq s
\]
\end{definition}

\begin{definition}[Continuous Dcpo]
\label{def:continuous-dcpo}
A dcpo $D$ is \textbf{continuous} if for every $x \in D$:
\[
x = \bigsqcup \{y \in D \mid y \ll x\}
\]
\end{definition}

\begin{theorem}[TKS Function Domains are Continuous]
\label{thm:function-continuous}
For continuous dcpos $D$ and $E$, the function space $[D \to E]$ of Scott continuous functions is a continuous dcpo.
\end{theorem}

%% ============================================================================
%% CHAPTER: EXTENDED DENOTATIONAL SEMANTICS
%% ============================================================================

\chapter{Extended Denotational Semantics}
\label{ch:extended-denotational}

\begin{tcolorbox}[colback=blue!5!white,colframe=blue!75!black,title=\vexpone]
This chapter extends the v7.0 denotational semantics with domain-theoretic foundations, logical relations, and full abstraction results.
\end{tcolorbox}

\section{Domain-Theoretic Semantic Function}

\begin{definition}[Pointed Semantic Domains]
\label{def:pointed-domains}
Extend the v7.0 semantic domains to pointed dcpos:

\textbf{Lifted Base Domains:}
\begin{align}
\sem{\mathsf{Int}}_\bot &= \mathbb{Z} \cup \{\bot\} \\
\sem{\mathsf{Bool}}_\bot &= \mathbb{B} \cup \{\bot\} \\
\sem{\mathsf{Element}[W]}_\bot &= \{Wn \mid n \in \{1,\ldots,10\}\} \cup \{\bot\}
\end{align}

\textbf{RPM Domain (refined):}
\[
\sem{\mathsf{RPM}[\tau]}_\bot = \mathcal{P}(\Acquisitions) \to (\sem{\tau}_\bot + \mathsf{Failure} + \{\bot\}) \times \mathcal{P}(\Acquisitions)
\]

\textbf{Continuous Function Domain:}
\[
\sem{\tau \to \sigma} = [\sem{\tau}_\bot \to \sem{\sigma}_\bot]_c
\]
where $[\cdot \to \cdot]_c$ denotes Scott continuous functions.
\end{definition}

\begin{definition}[Strict Semantic Function]
\label{def:strict-sem}
Define the \textbf{strict} semantic function $\sem{\cdot}^s$ that handles non-termination:

\textbf{Divergence:}
\[
\sem{\Omega}_\rho = \bot
\]
where $\Omega$ is a diverging term.

\textbf{Recursive Definitions:}
\[
\sem{\mathsf{fix}\; f}_\rho = \fix(\sem{f}_\rho)
\]
using the Kleene fixed-point.

\textbf{Strict Application:}
\[
\sem{e_1\; e_2}_\rho =
\begin{cases}
\bot & \text{if } \sem{e_1}_\rho = \bot \text{ or } \sem{e_2}_\rho = \bot \\
\sem{e_1}_\rho(\sem{e_2}_\rho) & \text{otherwise}
\end{cases}
\]
\end{definition}

\begin{theorem}[Semantic Function is Scott Continuous]
\label{thm:sem-scott}
For any well-typed term $\Gamma \types e : \tau$, the function:
\[
\sem{e} : \sem{\Gamma} \to \sem{\tau}_\bot
\]
is Scott continuous, where $\sem{\Gamma} = \prod_{x:\sigma \in \Gamma} \sem{\sigma}_\bot$.
\end{theorem}

\begin{proof}
By structural induction on typing derivations. Key cases:

\textbf{Case $\lambda x:\sigma. e$:} By IH, $\sem{e}$ is continuous in extended environments. Function abstraction preserves continuity.

\textbf{Case $e_1\; e_2$:} Composition of continuous functions is continuous. Application in continuous function spaces is continuous.

\textbf{Case $\mathsf{fix}\; f$:} By Theorem~\ref{thm:kleene-tks}, the fixed-point operator is continuous.

\textbf{Case $\noe{k}(e)$:} By Theorem~\ref{thm:noetic-scott}, noetic operators are continuous.
\end{proof}

\section{Logical Relations}

\begin{definition}[Logical Relation for TKS]
\label{def:logical-rel}
Define a family of relations $\mathcal{R}_\tau \subseteq \sem{\tau} \times \sem{\tau}$ indexed by types:

\textbf{Base Types:}
\[
\mathcal{R}_{\mathsf{Int}} = \{(n, n) \mid n \in \mathbb{Z}\} \cup \{(\bot, \bot)\}
\]

\textbf{Element Types:}
\[
\mathcal{R}_{\mathsf{Element}[W]} = \{(Wn, Wn) \mid n \in \{1,\ldots,10\}\} \cup \{(\bot, \bot)\}
\]

\textbf{Function Types:}
\[
(f, g) \in \mathcal{R}_{\tau \to \sigma} \iff \forall (v, v') \in \mathcal{R}_\tau.\, (f(v), g(v')) \in \mathcal{R}_\sigma
\]

\textbf{RPM Types:}
\begin{align}
(m, m') \in \mathcal{R}_{\mathsf{RPM}[\tau]} \iff &\forall \alpha.\, \text{if } m(\alpha) = (\mathsf{Success}(v), \alpha') \\
&\text{then } m'(\alpha) = (\mathsf{Success}(v'), \alpha') \text{ with } (v, v') \in \mathcal{R}_\tau
\end{align}
\end{definition}

\begin{lemma}[Fundamental Lemma]
\label{lem:fundamental}
If $\Gamma \types e : \tau$ and $(\rho, \rho') \in \mathcal{R}_\Gamma$ (pointwise), then:
\[
(\sem{e}_\rho, \sem{e}_{\rho'}) \in \mathcal{R}_\tau
\]
\end{lemma}

\begin{proof}
By induction on the typing derivation. The key insight is that all TKS operations preserve the logical relation.
\end{proof}

\section{Full Abstraction}

\begin{definition}[Observational Equivalence]
\label{def:obs-equiv}
Terms $e_1$ and $e_2$ are \textbf{observationally equivalent}, written $e_1 \approx_{obs} e_2$, if for all closing contexts $C[\cdot]$:
\[
C[e_1] \bigstep v \iff C[e_2] \bigstep v'
\]
\end{definition}

\begin{definition}[Denotational Equivalence]
\label{def:den-equiv}
Terms $e_1$ and $e_2$ are \textbf{denotationally equivalent}, written $e_1 \approx_{den} e_2$, if:
\[
\forall \rho.\, \sem{e_1}_\rho = \sem{e_2}_\rho
\]
\end{definition}

\begin{theorem}[Soundness (Adequacy)]
\label{thm:adequacy-full}
Denotational equivalence implies observational equivalence:
\[
e_1 \approx_{den} e_2 \implies e_1 \approx_{obs} e_2
\]
\end{theorem}

\begin{proof}
By the compositionality of $\sem{\cdot}$. If $\sem{e_1}_\rho = \sem{e_2}_\rho$ for all $\rho$, then in any context $C[\cdot]$:
\[
\sem{C[e_1]}_\rho = \sem{C}_{\rho'}[\sem{e_1}_\rho] = \sem{C}_{\rho'}[\sem{e_2}_\rho] = \sem{C[e_2]}_\rho
\]
Combined with Theorem~\ref{thm:semantic-soundness} from v7.0, the result follows.
\end{proof}

\begin{theorem}[Full Abstraction for TKS]
\label{thm:full-abstraction}
For the fragment of TKS without general recursion:
\[
e_1 \approx_{obs} e_2 \iff e_1 \approx_{den} e_2
\]
\end{theorem}

\begin{proof}
Soundness is Theorem~\ref{thm:adequacy-full}. For completeness, we use the logical relations technique.

Suppose $e_1 \not\approx_{den} e_2$. Then there exists $\rho$ with $\sem{e_1}_\rho \neq \sem{e_2}_\rho$.

Construct a distinguishing context $C[\cdot]$ as follows:
\begin{itemize}
  \item If the difference is at a base type, use equality testing
  \item If the difference is at function type, apply a suitable argument
  \item If the difference is at RPM type, run in a distinguishing acquisition state
\end{itemize}

The finite-type property of TKS base domains ensures this construction terminates.
\end{proof}

\begin{remark}
Full abstraction fails for TKS with general recursion, as is typical for domain-theoretic semantics. The counterexample involves parallel-or-like constructs that can be distinguished observationally but not denotationally in a sequential semantics.
\end{remark}

\section{Metric Semantics for Fractal Convergence}

\begin{definition}[Fractal Metric Space]
\label{def:fractal-metric}
Define a metric $d$ on fractals:
\[
d(\Frac_1, \Frac_2) =
\begin{cases}
0 & \text{if } \Frac_1 = \Frac_2 \\
2^{-n} & \text{where } n = \min\{i \mid \Frac_1[i] \neq \Frac_2[i]\} \\
1 & \text{if } \Frac_1 = \epsilon \text{ or } \Frac_2 = \epsilon
\end{cases}
\]
This is an ultrametric: $d(\Frac_1, \Frac_3) \leq \max(d(\Frac_1, \Frac_2), d(\Frac_2, \Frac_3))$.
\end{definition}

\begin{theorem}[Fractal Space Completeness]
\label{thm:fractal-complete}
The space $(\FracSet, d)$ is a complete metric space.
\end{theorem}

\begin{proof}
Let $(\Frac_n)_{n \in \mathbb{N}}$ be a Cauchy sequence. For each position $i$, the sequence $(\Frac_n[i])$ is eventually constant (by the Cauchy property with $\epsilon = 2^{-i-1}$). Define $\Frac^*[i] = \lim_n \Frac_n[i]$. Then $\Frac^* = \lim_n \Frac_n$.
\end{proof}

\begin{corollary}[Fractal Iteration Convergence]
\label{cor:fractal-convergence}
For a contractive fractal operator $F : \FracSet \to \FracSet$ (with Lipschitz constant $< 1$), the iteration sequence $\{F^n(\Frac_0)\}$ converges to a unique fixed point.
\end{corollary}

%% ============================================================================
%% CHAPTER: TOPOLOGICAL STRUCTURE OF NOETIC SPACE
%% ============================================================================

\chapter{Topological Structure of Noetic Space}
\label{ch:noetic-topology}

\begin{tcolorbox}[colback=blue!5!white,colframe=blue!75!black,title=\vnewone]
This chapter develops the topological and differential structure of noetic space, providing geometric intuition for noetic transformations.
\end{tcolorbox}

\section{Noetic Space as Topological Space}

\begin{definition}[Noetic Topological Space]
\label{def:noetic-topo}
The \textbf{noetic space} $\mathcal{N}$ is the set $\{\noe{0}, \noe{1}, \ldots, \noe{9}\}$ equipped with the discrete topology:
\[
\mathcal{O}(\mathcal{N}) = \mathcal{P}(\mathcal{N})
\]
\end{definition}

\begin{definition}[Extended Noetic Space]
\label{def:extended-noetic}
The \textbf{extended noetic space} $\bar{\mathcal{N}}$ includes limiting behavior:
\[
\bar{\mathcal{N}} = \mathcal{N} \cup \{\noe{\infty}\}
\]
with topology generated by:
\begin{itemize}
  \item All singletons $\{\noe{k}\}$ for $k \in \{0, \ldots, 9\}$
  \item Complements of finite sets: $\bar{\mathcal{N}} \setminus F$ for finite $F \subset \mathcal{N}$
\end{itemize}
This is the one-point compactification of $\mathcal{N}$.
\end{definition}

\begin{proposition}[Compactness of Extended Noetic Space]
\label{prop:compact-noetic}
$\bar{\mathcal{N}}$ is compact Hausdorff.
\end{proposition}

\section{World Manifold Structure}

\begin{definition}[World Manifold]
\label{def:world-manifold}
For each world $W$, define the \textbf{world manifold} $\mathcal{M}_W$ as the discrete 10-point space:
\[
\mathcal{M}_W = \{W1, W2, \ldots, W10\}
\]
equipped with a weighted graph structure encoding noetic connections.
\end{definition}

\begin{definition}[Inter-World Bundle]
\label{def:inter-world-bundle}
The \textbf{total world space} is the disjoint union:
\[
\mathcal{M} = \bigsqcup_{W \in \Worlds} \mathcal{M}_W = \mathcal{M}_A \sqcup \mathcal{M}_B \sqcup \mathcal{M}_C \sqcup \mathcal{M}_D
\]
with projection $\pi : \mathcal{M} \to \Worlds$ given by $\pi(Wn) = W$.

This forms a fiber bundle with discrete fibers.
\end{definition}

\begin{definition}[Descent Connection]
\label{def:descent-connection}
The \textbf{descent connection} $\nabla$ on the world bundle assigns to each descent path $\gamma : [0,1] \to \Worlds$ a parallel transport:
\[
\nabla_\gamma : \mathcal{M}_{\gamma(0)} \to \mathcal{M}_{\gamma(1)}
\]
For the elementary descent $A \to B$:
\[
\nabla_{A \to B}(An) = Bn
\]
\end{definition}

\begin{theorem}[ACBE as Parallel Transport]
\label{thm:acbe-transport}
The ACBE transformation is parallel transport along the canonical descent path $\gamma_{ACBE} : A \to B \to C \to D$:
\[
\mathrm{ACBE} = \nabla_{\gamma_{ACBE}} : \mathcal{M}_A \to \mathcal{M}_D
\]
\end{theorem}

\section{Sheaf-Theoretic Interpretation}

\begin{definition}[Presheaf of Elements]
\label{def:presheaf-elements}
Define the presheaf $\mathcal{E} : \mathbf{World}^{op} \to \mathbf{Set}$:
\begin{align}
\mathcal{E}(W) &= \{Wn \mid n \in \{1, \ldots, 10\}\} \\
\mathcal{E}(\iota_{W \to W'}) &= \text{(restriction is trivial for discrete fibers)}
\end{align}
\end{definition}

\begin{theorem}[Sheaf Condition]
\label{thm:sheaf-condition}
$\mathcal{E}$ is a sheaf on $\mathbf{World}$ with the natural topology.
\end{theorem}

\begin{definition}[Noetic Sheaf]
\label{def:noetic-sheaf}
The \textbf{noetic sheaf} $\mathcal{N}$ assigns to each world the available noetic operators:
\[
\mathcal{N}(W) = \{\noe{k} : \mathcal{E}(W) \to \mathcal{E}(W) \mid k \in \{0, \ldots, 9\}\}
\]
\end{definition}

\begin{proposition}[Noetic Sheaf Cohomology]
\label{prop:noetic-cohomology}
The first cohomology $H^1(\mathbf{World}, \mathcal{N})$ classifies twisted noetic structures---variants of TKS with non-trivial global monodromy around the descent cycle.
\end{proposition}

\section{Fractal Dimension Theory}

\begin{definition}[Fractal Dimension]
\label{def:fractal-dim}
For a fractal $\Frac = \langle k_1 : k_2 : \cdots \rangle$, define the \textbf{fractal dimension}:
\[
\dim(\Frac) = \lim_{n \to \infty} \frac{\log N_n(\Frac)}{\log n}
\]
where $N_n(\Frac)$ counts distinct subpaths of length $n$.
\end{definition}

\begin{theorem}[Eigenfractal Dimension]
\label{thm:eigenfractal-dim}
For a pure eigenfractal $\Frac_k = \langle k : k : k : \cdots \rangle$:
\[
\dim(\Frac_k) = 0
\]
For a maximally entropic fractal (all noetics equally likely):
\[
\dim(\Frac_{max}) = \log_{10}(10) = 1
\]
\end{theorem}

\begin{corollary}[Dimension Bounds]
\label{cor:dim-bounds}
For any fractal $\Frac$:
\[
0 \leq \dim(\Frac) \leq 1
\]
with equality at 0 iff $\Frac$ is eventually periodic with period 1.
\end{corollary}

%% ============================================================================
%% PART III: ADVANCED RPM STRUCTURES (v7.1 EXTENSION)
%% ============================================================================
\part{Advanced RPM Theory}

\chapter{Multi-Goal RPM Algebra}
\label{ch:multi-goal-rpm}

\begin{tcolorbox}[colback=blue!5!white,colframe=blue!75!black,title=\vnewone]
This chapter extends RPM to handle multiple simultaneous goals, products, and coproducts of prerequisite structures.
\end{tcolorbox}

\section{RPM Products}

\begin{definition}[RPM Product]
\label{def:rpm-product}
For RPM computations $m_1 : \RPM[\tau_1]$ and $m_2 : \RPM[\tau_2]$, the \textbf{product} is:
\[
m_1 \times_\RPM m_2 : \RPM[\tau_1 \times \tau_2]
\]
defined by:
\[
(m_1 \times_\RPM m_2)(\alpha) =
\begin{cases}
(\mathsf{Success}(v_1, v_2), \alpha'') & \text{if } m_1(\alpha) = (\mathsf{Success}(v_1), \alpha') \\
                                        & \text{and } m_2(\alpha') = (\mathsf{Success}(v_2), \alpha'') \\
(\mathsf{Failure}(p), \alpha') & \text{if } m_1(\alpha) = (\mathsf{Failure}(p), \alpha') \\
(\mathsf{Failure}(p), \alpha'') & \text{if } m_1(\alpha) = (\mathsf{Success}(\cdot), \alpha') \\
                                 & \text{and } m_2(\alpha') = (\mathsf{Failure}(p), \alpha'')
\end{cases}
\]
\end{definition}

\begin{proposition}[Product Associativity]
\label{prop:product-assoc}
RPM products are associative up to isomorphism:
\[
(m_1 \times_\RPM m_2) \times_\RPM m_3 \cong m_1 \times_\RPM (m_2 \times_\RPM m_3)
\]
\end{proposition}

\begin{definition}[Parallel RPM]
\label{def:parallel-rpm}
The \textbf{parallel product} runs both computations independently:
\[
m_1 \otimes m_2 : \RPM[\tau_1 \times \tau_2]
\]
\[
(m_1 \otimes m_2)(\alpha) =
\begin{cases}
(\mathsf{Success}(v_1, v_2), \alpha_1 \cup \alpha_2) & \text{if both succeed} \\
(\mathsf{Failure}(p_1 \lor p_2), \alpha) & \text{if either fails}
\end{cases}
\]
where $m_i(\alpha) = (\mathsf{Success}(v_i), \alpha_i)$ or $(\mathsf{Failure}(p_i), \_)$.
\end{definition}

\section{RPM Coproducts}

\begin{definition}[RPM Coproduct]
\label{def:rpm-coproduct}
For RPM computations $m_1 : \RPM[\tau]$ and $m_2 : \RPM[\tau]$, the \textbf{coproduct} (choice):
\[
m_1 +_\RPM m_2 : \RPM[\tau]
\]
defined by:
\[
(m_1 +_\RPM m_2)(\alpha) =
\begin{cases}
m_1(\alpha) & \text{if } m_1(\alpha) = (\mathsf{Success}(\_), \_) \\
m_2(\alpha) & \text{otherwise}
\end{cases}
\]
\end{definition}

\begin{theorem}[RPM as Distributive Category]
\label{thm:rpm-distributive}
The category of RPM computations with $\times_\RPM$ and $+_\RPM$ forms a distributive category:
\[
m \times_\RPM (n_1 +_\RPM n_2) \cong (m \times_\RPM n_1) +_\RPM (m \times_\RPM n_2)
\]
\end{theorem}

\section{Multi-Goal Pursuit}

\begin{definition}[Goal Set]
\label{def:goal-set}
A \textbf{goal set} is a finite collection $G = \{g_1, \ldots, g_n\}$ where each $g_i : \RPM[\tau_i]$.
\end{definition}

\begin{definition}[Goal Conjunction]
\label{def:goal-conj}
The \textbf{conjunction} of goals:
\[
\bigwedge G = g_1 \times_\RPM g_2 \times_\RPM \cdots \times_\RPM g_n : \RPM[\tau_1 \times \cdots \times \tau_n]
\]
succeeds only if all goals succeed.
\end{definition}

\begin{definition}[Goal Disjunction]
\label{def:goal-disj}
The \textbf{disjunction} of goals (where all $\tau_i = \tau$):
\[
\bigvee G = g_1 +_\RPM g_2 +_\RPM \cdots +_\RPM g_n : \RPM[\tau]
\]
succeeds if any goal succeeds.
\end{definition}

\begin{theorem}[Multi-Goal Soundness]
\label{thm:multi-goal-sound}
For goal set $G = \{g_1, \ldots, g_n\}$ and acquisition state $\alpha$:
\begin{enumerate}
  \item $(\bigwedge G)(\alpha) = \mathsf{Success}(\bar{v})$ iff $\forall i.\, g_i$ eventually succeeds
  \item $(\bigvee G)(\alpha) = \mathsf{Success}(v)$ iff $\exists i.\, g_i(\alpha) = \mathsf{Success}(v, \_)$
\end{enumerate}
\end{theorem}

%% ============================================================================
%% VERSION HISTORY
%% ============================================================================

\chapter{Version History}
\label{app:versions}

\begin{description}
\item[v7.1 (Current)] Advanced Formal Semantics Release
  \begin{itemize}
  \item Natural transformations between worlds
  \item Domain-theoretic foundations (dcpos, Scott topology)
  \item Extended denotational semantics with full abstraction
  \item Topological structure of noetic space
  \item Multi-goal RPM algebra (products, coproducts)
  \end{itemize}

\item[v7.0] Academic Expansion Release
  \begin{itemize}
  \item Complete Noetic Fractal Calculus with type theory and eigenfractals
  \item RPM Kleisli category and effect system
  \item Complete dynamic semantics for all constructs
  \item Full Progress and Preservation proofs
  \item Concurrency model prototype
  \item Quantum Noetic algebra prototype
  \item 20+ worked examples
  \end{itemize}

\item[v6.1] Dynamic Semantics Release
  \begin{itemize}
  \item RPM formalized as monad with proofs
  \item Basic Noetic Fractal Calculus
  \item Integrated evaluation rules
  \item Extended examples
  \end{itemize}

\item[v6.0] Academic Release

\item[v5.0] Mathematical Framework

\item[v3.2.1] Initial Formalization
\end{description}

%% ============================================================================
%% BACKMATTER
%% ============================================================================
\backmatter

\begin{thebibliography}{99}
\bibitem{kybalion} Three Initiates. \textit{The Kybalion}. Yogi Publication Society, 1908.
\bibitem{moggi} E. Moggi. ``Notions of computation and monads.'' \textit{Information and Computation}, 93(1):55--92, 1991.
\bibitem{wadler} P. Wadler. ``Monads for functional programming.'' \textit{Advanced Functional Programming}, LNCS 925:24--52, 1995.
\bibitem{pierce} B. C. Pierce. \textit{Types and Programming Languages}. MIT Press, 2002.
\bibitem{plotkin} G. D. Plotkin. ``A structural approach to operational semantics.'' DAIMI FN-19, 1981.
\bibitem{mac-lane} S. Mac Lane. \textit{Categories for the Working Mathematician}. Springer, 1971.
\bibitem{plotkin-power} G. Plotkin and J. Power. ``Algebraic operations and generic effects.'' \textit{Applied Categorical Structures}, 11(1):69--94, 2003.
\bibitem{nielsen-chuang} M. Nielsen and I. Chuang. \textit{Quantum Computation and Quantum Information}. Cambridge University Press, 2000.
\bibitem{abramsky-jung} S. Abramsky and A. Jung. ``Domain Theory.'' \textit{Handbook of Logic in Computer Science}, Vol. 3, Oxford, 1994.
\bibitem{scott} D. S. Scott. ``Data types as lattices.'' \textit{SIAM Journal on Computing}, 5(3):522--587, 1976.
\bibitem{milner} R. Milner. ``Fully abstract models of typed λ-calculi.'' \textit{TCS}, 4(1):1--22, 1977.
\end{thebibliography}

\end{document}
